% ==============================================================================
% Fichier : chapitres/02_security_risk_management.tex
% Livre II : Security and Risk Management
% ==============================================================================

\chapter{Security and Risk Management}
\label{chap:risk_management}

\begin{boxobjectifs}
\textbf{Objectifs du Livre~II :} Donner les éléments de compréhension liés à la gestion des risques, domaine majeur de la sécurité informatique. À l'issue, l'étudiant saura élaborer une politique de sécurité, mener une analyse des risques, identifier les cadres légaux applicables et distinguer les types d'attaquants.
\end{boxobjectifs}

% ------------------------------------------------------------------------------
\section{Gouvernance de la Sécurité de l'Information}
% ------------------------------------------------------------------------------

\subsection{Élaboration des Politiques de Sécurité}

Une politique de sécurité est un document officiel qui définit les règles et les directives adoptées par une organisation pour gérer et protéger ses actifs informationnels. Elle doit être :

\begin{itemize}
    \item \textbf{Alignée sur les objectifs stratégiques} de l'organisation.
    \item \textbf{Approuvée par la direction} pour lui conférer son autorité.
    \item \textbf{Communiquée} à l'ensemble du personnel.
    \item \textbf{Régulièrement révisée} pour tenir compte de l'évolution des menaces.
\end{itemize}

\subsection{Procédures de Sécurité}

Les procédures de sécurité déclinent opérationnellement les politiques en étapes concrètes. Elles doivent :
\begin{enumerate}
    \item Être développées en concertation avec les parties prenantes.
    \item Faire l'objet de tests et de validations.
    \item Être communiquées et comprises par tout le personnel concerné.
    \item Être soigneusement documentées et archivées.
    \item Être évaluées régulièrement pour identifier les lacunes.
\end{enumerate}

\subsection{Rôles et Responsabilités}

\begin{table}[H]
\centering
\begin{tabular}{p{3.5cm}p{9cm}}
    \toprule
    \textbf{Acteur} & \textbf{Responsabilités} \\
    \midrule
    \textbf{DSI} (Directeur des SI) & Supervision globale du SI, définition de la stratégie sécurité, alignement sur les objectifs métier, allocation des ressources. \\
    \midrule
    \textbf{RSSI} (Resp. Sécurité SI) & Élaboration et mise en œuvre des politiques et procédures, évaluation des risques, supervision des incidents, sensibilisation du personnel. \\
    \midrule
    \textbf{Responsables métier} & Application des politiques dans leur périmètre, remontée des incidents. \\
    \midrule
    \textbf{Auditeurs} & Évaluation indépendante de la conformité et de l'efficacité des contrôles. \\
    \bottomrule
\end{tabular}
\caption{Principaux rôles dans la gouvernance de la sécurité de l'information}
\end{table}

% ------------------------------------------------------------------------------
\section{Contrôle d'Accès}
% ------------------------------------------------------------------------------

Les contrôles d'accès régulent l'accès aux systèmes d'information et se répartissent en deux grandes catégories :

\subsection{Contrôles d'Accès Physique}
\begin{itemize}
    \item Barrières physiques (portiques, tourniquets, sas de sécurité).
    \item Systèmes \og{}\textit{mantrap}\fg{} (accès d'une personne à la fois).
    \item Badges d'accès (RFID, magnétiques).
    \item Surveillance vidéo (CCTV).
    \item Gardiens et rondes de sécurité.
\end{itemize}

\subsection{Contrôles d'Accès Technico-logiques}
\begin{itemize}
    \item Listes de contrôle d'accès (ACL).
    \item Pares-feux filtrant le trafic réseau.
    \item Systèmes IDS/IPS surveillant les activités suspectes.
    \item Authentification multifacteur (MFA).
    \item Chiffrement des données en transit et au repos.
\end{itemize}

\begin{boxinfo}
La combinaison des contrôles physiques et logiques est essentielle pour garantir une sécurité globale. Ils protègent à la fois les infrastructures matérielles et les données numériques.
\end{boxinfo}

% ------------------------------------------------------------------------------
\section{Analyse des Risques}
% ------------------------------------------------------------------------------

\begin{center}
\begin{tikzpicture}[
    node distance=1.5cm,
    box/.style={rectangle, rounded corners, draw=CouleurPrimaire, fill=CouleurPrimaire!10,
                text width=3cm, align=center, minimum height=1cm, font=\small\bfseries}
]
    \node[box] (id)    {Identification\\des risques};
    \node[box] (eval)  [right=of id]   {Évaluation\\des risques};
    \node[box] (trait) [right=of eval] {Traitement\\des risques};
    \node[box] (surv)  [below=of trait] {Surveillance\\et révision};
    \node[box] (comm)  [left=of surv]  {Communication\\des risques};

    \draw[->, thick, CouleurPrimaire] (id) -- (eval);
    \draw[->, thick, CouleurPrimaire] (eval) -- (trait);
    \draw[->, thick, CouleurPrimaire] (trait) -- (surv);
    \draw[->, thick, CouleurPrimaire] (surv) -- (comm);
    \draw[->, thick, CouleurPrimaire] (comm) -- (id);
\end{tikzpicture}
\captionof{figure}{Cycle de gestion des risques}
\end{center}

\subsection{Identification des Risques}

Techniques courantes d'identification :
\begin{itemize}
    \item \textbf{Analyse documentaire :} Examen des politiques, rapports d'incidents et audits passés.
    \item \textbf{Entretiens et ateliers :} Recueil des perceptions des parties prenantes.
    \item \textbf{Analyse de flux :} Cartographie des traitements et des échanges de données.
    \item \textbf{Scans de vulnérabilités :} Outils automatisés identifiant les failles techniques.
\end{itemize}

\subsection{Évaluation des Risques}

\[
\text{Niveau de Risque} = \text{Probabilité d'occurrence} \times \text{Impact potentiel}
\]

\begin{table}[H]
\centering
\small
\begin{tabular}{lccc}
    \toprule
    & \textbf{Impact Faible} & \textbf{Impact Moyen} & \textbf{Impact Élevé} \\
    \midrule
    \textbf{Probabilité Élevée}  & Moyen  & Élevé    & Critique \\
    \textbf{Probabilité Moyenne} & Faible & Moyen    & Élevé    \\
    \textbf{Probabilité Faible}  & Faible & Faible   & Moyen    \\
    \bottomrule
\end{tabular}
\caption{Matrice de criticité des risques}
\end{table}

\subsection{Traitement du Risque}

\begin{itemize}
    \item \textbf{Réduction (Mitigation) :} Mise en place de contrôles pour diminuer la probabilité ou l'impact.
    \item \textbf{Transfert :} Externalisation du risque (assurance, sous-traitance contractuelle).
    \item \textbf{Acceptation :} Décision consciente d'assumer le risque résiduel.
    \item \textbf{Évitement :} Suppression de l'activité ou de l'actif générateur de risque.
\end{itemize}

% ------------------------------------------------------------------------------
\section{Lois et Régulations}
% ------------------------------------------------------------------------------

\subsection{Lois Internationales}
\begin{itemize}
    \item \textbf{RGPD} (Règlement Général sur la Protection des Données, UE 2018) : Protection des données personnelles des citoyens européens.
    \item \textbf{CFAA} (Computer Fraud and Abuse Act, USA) : Cadre pénal américain contre la fraude informatique.
    \item \textbf{Convention de Budapest} (2001) : Premier traité international sur la cybercriminalité.
\end{itemize}

\subsection{Lois Nationales (Cameroun)}
\begin{itemize}
    \item \textbf{Loi n°~2010/012} relative à la cybersécurité et à la cybercriminalité.
    \item Rôle de l'\textbf{ANTIC} (Agence Nationale des Technologies de l'Information et de la Communication) dans la régulation et la supervision de la sécurité des SI.
\end{itemize}

% ------------------------------------------------------------------------------
\section{Sécurité de Tierces Parties}
% ------------------------------------------------------------------------------

\subsection{Responsabilités envers les Tierces Parties}
L'organisation reste responsable de la sécurité des données, même lorsque leur traitement est confié à un tiers (prestataire, fournisseur Cloud, sous-traitant). La due diligence s'impose.

\subsection{Cadre de Gestion des Risques Tiers}
\begin{itemize}
    \item Évaluation de la maturité sécurité du prestataire (questionnaires, audits).
    \item Clauses contractuelles de sécurité (SLA, confidentialité, notification d'incidents).
    \item Surveillance continue des accès du prestataire au SI.
    \item Plan de sortie prévu en cas de défaillance du tiers.
\end{itemize}

% ------------------------------------------------------------------------------
\section{Éthique et Déontologie}
% ------------------------------------------------------------------------------

\begin{boxalert}
\textbf{Code d'éthique du professionnel de sécurité (ISACA/ISC²) :}
\begin{itemize}
    \item Protéger la société, les infrastructures communes et la cybersphère.
    \item Agir honorablement, honnêtement, justement, responsablement et légalement.
    \item Fournir des services diligents et compétents aux clients et aux employeurs.
    \item Promouvoir et protéger la profession.
\end{itemize}
\end{boxalert}

% ------------------------------------------------------------------------------
\section{Types d'Attaquants}
% ------------------------------------------------------------------------------

\begin{table}[H]
\centering
\small
\begin{tabular}{p{3cm}p{4cm}p{5cm}}
    \toprule
    \textbf{Type} & \textbf{Profil} & \textbf{Motivation} \\
    \midrule
    \textbf{Script Kiddie}    & Néophyte utilisant des outils prêts-à-l'emploi & Notoriété, curiosité \\
    \textbf{Hacker Black Hat} & Expert malveillant & Gain financier, espionnage \\
    \textbf{Hacker White Hat} & Expert éthique (pentest autorisé) & Améliorer la sécurité \\
    \textbf{Hacker Grey Hat}  & Entre les deux & Variable \\
    \textbf{Insider Threat}   & Employé malveillant ou négligent & Vengeance, gain, erreur humaine \\
    \textbf{Hacktiviste}      & Groupe idéologique & Protestation politique \\
    \textbf{État-nation}      & Service de renseignement & Espionnage, sabotage stratégique \\
    \bottomrule
\end{tabular}
\caption{Typologie des attaquants}
\end{table}

% ------------------------------------------------------------------------------
\section{Évaluation des Acquis --- Livre~II}
% ------------------------------------------------------------------------------

\begin{boxexo}
\textbf{QCM :}
\begin{enumerate}
    \item La formule $\text{Risque} = \text{Probabilité} \times \text{Impact}$ est utilisée dans quelle étape de la gestion des risques ? \textbf{(a) Identification \quad (b) Évaluation \quad (c) Traitement \quad (d) Surveillance}
    \item Quel est le rôle principal du RSSI ? \textbf{(a) Gérer le budget IT \quad (b) Élaborer et mettre en œuvre la politique sécurité \quad (c) Administrer les serveurs \quad (d) Développer les applications}
    \item Le transfert de risque consiste à : \textbf{(a) Supprimer le risque \quad (b) Le confier à un tiers (assurance) \quad (c) L'ignorer \quad (d) Réduire son impact}
    \item Un attaquant utilisant des outils trouvés sur Internet sans comprendre leur fonctionnement est appelé : \textbf{(a) Hacker Black Hat \quad (b) Insider Threat \quad (c) Script Kiddie \quad (d) Hacktiviste}
\end{enumerate}

\textbf{Travail Pratique :}\\
Réaliser une analyse de risque simplifiée pour un système d'information fictif en utilisant la matrice de criticité. Identifier trois risques, les évaluer et proposer un traitement pour chacun.
\end{boxexo}
